\documentclass[11pt,a4paper]{article}
\usepackage[T1]{fontenc}
\usepackage[utf8]{inputenc}
\usepackage{lmodern}
\usepackage{amsmath,amssymb}
\usepackage[margin=25mm]{geometry}
\usepackage{longtable,booktabs,array,calc}
\usepackage[hidelinks]{hyperref}
\hypersetup{pdftitle={Brownian free motion: the velocitas ultima denied, one action constant, and the floor},pdfauthor={navstokgap research working notes}}
\newcommand{\tightlist}{\setlength{\itemsep}{0pt}\setlength{\parskip}{0pt}}
\setlength{\emergencystretch}{3em}
\setcounter{secnumdepth}{0}
\begin{document}
\hypertarget{brownian-free-motion-the-velocitas-ultima-denied-one-action-constant-and-the-floor}{%
\section{Brownian free motion: the velocitas ultima denied, one action
constant, and the
floor}\label{brownian-free-motion-the-velocitas-ultima-denied-one-action-constant-and-the-floor}}

\textbf{Result, 2026-09-28 (Claude; written derivation, to be
refereed).} A third route for the
\href{principia-fifth-postulate.md}{fifth-postulate programme}, beside
the deformation route (Theorems A, B) and the Gaussian state route
(Theorem B\('\)): the stochastic one. Let the free motion of a body of
mass \(m\) be a Markov process with continuous paths, homogeneous in
space and time and isotropic (premises P1--P3 of §1).

\begin{itemize}
\tightlist
\item
  \textbf{Theorem S1 (the law and the denial).} The free motion is
  Brownian motion with drift, with variance rate \(D(m)\ge0\) per
  coordinate. For \(D=0\) it is Newton's inertial motion. For \(D>0\)
  the paths are almost surely nowhere differentiable, so the
  \emph{velocitas ultima} that Newton's scholium calls ``certus \&
  definitus'' exists only as a mean velocity.
\item
  \textbf{Theorem S2 (one constant).} If the law depends on the mass
  alone (P4) and the centre of mass of independent bodies obeys the law
  of their total mass (P5), then \(\kappa=mD\) is one universal
  constant, with the dimension of an action; \(\kappa=0\) is admissible.
\item
  \textbf{Theorem S3 (the floor).} For \(\kappa>0\), any record that
  decides from the path, with pinned ends, between the inertial chord
  and the constant-force parabola over a cell errs at equal priors with
  probability at least \(\Phi(-\sqrt{K_\tau/2\kappa})\), attained by the
  optimal test, where \(K_\tau=F^2\tau^3/(24m)\) is the Galileo action
  of the cell. Confidence \(1-\epsilon\) needs
  \(K_\tau\ge2z_{1-\epsilon}^2\kappa\); with \(\kappa=\hbar\) this is
  the mark mesh of the \href{planck-gap-paper.md}{Planck paper}, eq.
  (2).
\item
  \textbf{Proposition S4 (jumps allowed: the leap).} Dropping
  continuity, P2--P5 force the Lévy exponent of a body of mass \(m\) to
  be \(m\,\Gamma(|k|/m)\) for one universal function \(\Gamma\). If the
  only universal constant is an action \(\kappa\), dimensional analysis
  forces \(\Gamma(p)=C\kappa p^2\), and continuity follows. With a speed
  \(c\) as well, the relativistic kinetic energy divided by \(\kappa\),
  \((\sqrt{m^2c^4+c^2\kappa^2k^2}-mc^2)/\kappa\), is of this form and
  generates a process with jumps.
\item
  \textbf{Independence (Nelson).} With diffusion coefficient
  \(\nu=\kappa/2m\), Nelson's stochastic mechanics obeys Newton's second
  law in a mean sense and is equivalent to the Schrödinger equation with
  \(\hbar=\kappa\)
  (\href{https://doi.org/10.1103/PhysRev.150.1079}{Nelson 1966},
  abstract; precursor \href{https://doi.org/10.1007/BF01338578}{Fényes
  1952}, metadata). So every \(\kappa\ge0\) is consistent with the Laws,
  the analogue of Theorem A.
\end{itemize}

The route carries the four parts of the fifth-postulate goal: the
statement (the existence of the ultimate velocity), its independence,
exactly one added action constant, and a floor on recording the
inertial--parabola difference, an area Newton takes to zero. Its limit
is the same as in the other routes: \(\kappa>0\) is the negated
postulate, assumed, and the zero branch stays admissible. Theorems S1
and S2 are elementary; the centre-of-mass argument of S2 is likely known
in stochastic mechanics, where \(\nu\propto1/m\) is postulated, but it
has not been located in print here.

\hypertarget{premises}{%
\subsection{1. Premises}\label{premises}}

\begin{itemize}
\tightlist
\item
  \textbf{P1 (continuity).} The free motion is a Markov process with
  continuous paths. Continuity of motion is Aristotle's (\emph{Physics}
  VI, cited by book only); the leap of al-Nazzam, already in the Planck
  paper's classics, would be a jump, which P1 excludes.
\item
  \textbf{P2 (homogeneity).} Stationary independent increments: no
  preferred place or time. A Galilean boost adds a constant drift and
  preserves P2.
\item
  \textbf{P3 (isotropy).} No preferred direction.
\item
  \textbf{P4 (mass alone).} A body's free law depends only on its mass,
  as all bodies fall alike.
\item
  \textbf{P5 (composition).} For independent bodies the centre of mass
  obeys the free law of the total mass.
\end{itemize}

\hypertarget{theorem-s1}{%
\subsection{2. Theorem S1}\label{theorem-s1}}

A process in \(\mathbb R^3\) with stationary independent increments and
continuous paths is a Brownian motion with drift,
\(X_t=x_0+vt+\Sigma^{1/2}W_t\) (the Lévy--Itô decomposition has no jump
part when the paths are continuous; standard). Isotropy gives
\(\Sigma=D\,I\), and \(D\ge0\) can depend only on the body. For \(D>0\),
Brownian paths are almost surely nowhere differentiable
(\href{https://doi.org/10.1007/BF01474606}{Paley, Wiener and Zygmund
1933}, metadata), and they have Hausdorff dimension 2, the dimension
found for quantum paths by \href{https://doi.org/10.1119/1.12657}{Abbott
and Wise (1981)} (metadata). The difference quotient \((X_{t+h}-X_t)/h\)
has no limit; its conditional mean does,
\(\lim_{h\downarrow0}E[X_{t+h}-X_t\mid\mathcal F_t]/h=v\), which is
Nelson's mean forward velocity. \(\square\)

In the words of the scholium closing Book I, Section I (quoted from the
1687 text in §1 of the
\href{principia-fifth-postulate.md}{fifth-postulate note}): ``Extat
limes quem velocitas in fine motus attingere potest, non autem
transgredi. Hæc est velocitas ultima. {[}\ldots{]} Cumq; hic limes sit
certus \& definitus''. For \(D>0\) the limit of the ratio exists for the
mean and fails for the path. This is the denial of joint determinacy in
the form Newton wrote it: the place at an instant is sharp, while the
velocity at that instant has no pathwise value.

\hypertarget{theorem-s2}{%
\subsection{3. Theorem S2}\label{theorem-s2}}

\emph{Proof.} Take independent bodies with masses \(m_i\), drifts
\(v_i\) and variance rates \(D(m_i)\), and total mass \(M=\sum m_i\).
The centre of mass \(X=\sum m_ix_i/M\) has the constant drift
\(\sum m_iv_i/M\) (momentum conservation) and the variance rate
\(\sum m_i^2D(m_i)/M^2\). By P5 this equals \(D(M)\), so the function
\(g(m)=m^2D(m)\) satisfies \(g(m_1+m_2)=g(m_1)+g(m_2)\) for all
\(m_1,m_2>0\). Since \(g\ge0\), \(g\) is nondecreasing, and Cauchy's
equation gives \(g(m)=\kappa m\). Hence \(D(m)=\kappa/m\) and
\(mD=\kappa\) is the same for every body. Its dimension is mass times
length\(^2\)/time, an action. \(\square\)

The argument is the stochastic counterpart of the
\href{rotation-composition-universality.md}{rotation-composition note},
where interaction forces a common constant in the composition of
rotations; P4 plays the role of the equivalence principle, and P5 the
role of the centre-of-mass law. Neither argument excludes \(\kappa=0\).

\hypertarget{theorem-s3}{%
\subsection{4. Theorem S3}\label{theorem-s3}}

Let a constant force \(F\) act, so the body's path is the Newtonian
parabola plus \(\sqrt{D}\,W_t\). Pin the ends of a cell of duration
\(\tau\). The two laws to be told apart, with and without the force, are
Brownian bridges of variance rate \(D=\kappa/m\) around the parabola and
around the chord. Their difference is the shift \(\delta x\) (chord
minus parabola), with Cameron--Martin norm
\(\|\delta x\|^2=D^{-1}\int_0^\tau\dot{\delta x}^2dt=(m/\kappa)\int\dot{\delta x}^2=2K_\tau/\kappa\).
The Neyman--Pearson test with equal priors errs with probability
\(\Phi(-\frac12\|\delta x\|)=\Phi(-\sqrt{K_\tau/2\kappa})\), and no test
does better. This is Proposition 6 of the
\href{cut-measure-newton.md}{cut-measure note} with \(\hbar\) replaced
by \(\kappa\); its Theorem 2 distributes the evidence \(K_\tau/\kappa\)
over any sequence of cuts, each cut carrying the Kullback--Leibler share
\(3s(1-s)K_L/\kappa\). \(\square\)

With \(\kappa=\hbar\), which is the normalization of the Euclidean free
measure \(e^{-S/\hbar}\) and of Nelson's \(\nu=\hbar/2m\), the threshold
is \(\tau\ge(48z_{1-\epsilon}^2m\hbar/F^2)^{1/3}\), the Planck paper's
eq. (2).

\hypertarget{where-the-route-stands}{%
\subsection{5. Where the route stands}\label{where-the-route-stands}}

\begin{longtable}[]{@{}
  >{\raggedright\arraybackslash}p{(\columnwidth - 6\tabcolsep) * \real{0.2500}}
  >{\raggedright\arraybackslash}p{(\columnwidth - 6\tabcolsep) * \real{0.2500}}
  >{\raggedright\arraybackslash}p{(\columnwidth - 6\tabcolsep) * \real{0.2500}}
  >{\raggedright\arraybackslash}p{(\columnwidth - 6\tabcolsep) * \real{0.2500}}@{}}
\toprule\noalign{}
\begin{minipage}[b]{\linewidth}\raggedright
part of the goal
\end{minipage} & \begin{minipage}[b]{\linewidth}\raggedright
deformation route
\end{minipage} & \begin{minipage}[b]{\linewidth}\raggedright
state route
\end{minipage} & \begin{minipage}[b]{\linewidth}\raggedright
stochastic route
\end{minipage} \\
\midrule\noalign{}
\endhead
\bottomrule\noalign{}
\endlastfoot
statement changed & commutativity (Thm A) & sharp states (Thm B\('\)) &
\emph{velocitas ultima} exists pathwise (S1) \\
independence from the Laws & Moyal (Thm A) & Gaussian closure & Nelson
1966 \\
exactly one constant & Gutt, under covariance (Thm B) & Thm B\('\) &
composition, P4--P5 (S2) \\
floor on the comparison & Thms C, E & Thm C & Cameron--Martin (S3) \\
\end{longtable}

The premise that remains is the same in all three columns: why the
constant is positive. In the stochastic route it has the plainest form,
whether free paths have an instantaneous velocity, and Newton's scholium
asserts that they do. Two further limits: P4 and P5 are premises about
bodies (an equivalence principle for fluctuations and a centre-of-mass
law), and S3 is proved for additive position noise, the law of S1 with a
force added; in Nelson's full dynamics the drift depends on the state,
and the corresponding bound is not claimed here. Allowing jumps (the
leap) is treated in §5b.

\hypertarget{b.-jumps-allowed-proposition-s4}{%
\subsection{5b. Jumps allowed (Proposition
S4)}\label{b.-jumps-allowed-proposition-s4}}

Drop P1 and keep P2--P5. The free law of a body of mass \(m\) is then an
isotropic Lévy process; take it symmetric, so its exponent,
\(E\,e^{ik\cdot(X_t-X_0)}=e^{-t\,\phi(m,|k|)}\), is real and
nonnegative. The centre of mass of independent bodies has exponent
\(\sum_i\phi(m_i,m_i|k|/M)\), and P5 requires
\(\phi(M,|k|)=\sum_i\phi(m_i,m_i|k|/M)\). Put \(|k|=Mp\) and
\(G(m,p)=\phi(m,mp)\): then \(G(m_1+m_2,p)=G(m_1,p)+G(m_2,p)\), and
\(G\ge0\) gives \(G(m,p)=m\,\Gamma(p)\) with \(\Gamma(p)=G(1,p)\). Hence

\[\phi(m,|k|)=m\,\Gamma(|k|/m),\qquad\nu_m(A)=m\,\nu(mA),\]

where \(\nu\) is the Lévy measure of \(\Gamma\): a body of mass \(m\)
jumps with the universal jumps shrunk by \(1/m\) and their rate
multiplied by \(m\). Brownian motion is \(\Gamma(p)=\kappa p^2/2\),
giving \(D=\kappa/m\) as in S2, and the symmetric \(\alpha\)-stable laws
are \(\Gamma(p)=\kappa_\alpha p^\alpha\), with scale
\(\kappa_\alpha m^{1-\alpha}\). \(\square\)

\emph{Dimensional analysis.} \(\phi\) is an inverse time and \(p=|k|/m\)
has dimension \(({\rm mass}\cdot{\rm length})^{-1}\), so \(\Gamma\) has
dimension \(({\rm mass}\cdot{\rm time})^{-1}\). If the only universal
constant is one action \(\kappa\), the unique monomial of that dimension
is \(\kappa p^2\), and no dimensionless combination of \(\kappa\) and
\(p\) exists; so \(\Gamma(p)=C\kappa p^2\) and the law is Brownian.
Continuity is then a consequence: a jump needs a second constant. With a
speed \(c\) available, \(\kappa p/c\) is dimensionless and
\(\Gamma(p)=\kappa p^2f(\kappa p/c)\) is allowed. The choice
\(\Gamma(p)=(\sqrt{c^4+c^2\kappa^2p^2}-c^2)/\kappa\) gives
\(\phi=m\,\Gamma(|k|/m)=(\sqrt{m^2c^4+c^2\kappa^2k^2}-mc^2)/\kappa\),
the relativistic kinetic energy divided by \(\kappa\), which reduces to
\(\kappa k^2/2m\) for \(\kappa|k|\ll mc\); it is a Bernstein function of
\(k^2\), and its process has jumps
(\href{https://doi.org/10.1016/0022-1236(90)90049-Q}{Carmona, Masters
and Simon 1990}, metadata). So in this reading the leap is what a speed
limit readmits, and Aristotle's continuity is what a single action
constant enforces. This is a structural remark: which constants the free
law may contain is itself a premise.

\hypertarget{consequence-for-state}{%
\subsection{6. Consequence for STATE}\label{consequence-for-state}}

Newton necessity gains a third route in which the negated postulate is
Newton's own sentence on the ultimate velocity: continuity, homogeneity
and isotropy give Brownian free motion, composition gives one action
constant, and the Cameron--Martin theorem gives the floor, equal to the
Planck paper's mark mesh at \(\kappa=\hbar\). Positivity of \(\kappa\)
remains a premise.
\end{document}
